\chapter{Kontinuierliche Zufallsvariablen}

\section{Ziel dieses Kapitels}

Im vorherigen Kapitel wurden diskrete Zufallsvariablen (discrete random variables) betrachtet.
Eine diskrete Zufallsvariable kann nur endlich viele oder abzählbar viele Werte annehmen.
Viele Größen in Anwendungen sind aber nicht sinnvoll diskret, sondern kontinuierlich (continuous).

Beispiele sind:

\begin{itemize}
    \item Körpergröße,
    \item Temperatur,
    \item Zeit,
    \item Geschwindigkeit,
    \item Position auf einer Linie,
    \item Messwerte mit Rauschen (noise),
    \item Pixelintensitäten, wenn sie als reelle Werte modelliert werden.
\end{itemize}

In diesem Kapitel lernen wir, wie man solche Größen mit Wahrscheinlichkeitsdichten (probability density functions) beschreibt.

\begin{conceptbox}
Bei diskreten Zufallsvariablen ordnet man einzelnen Werten Wahrscheinlichkeiten zu.

Bei kontinuierlichen Zufallsvariablen ordnet man einzelnen Werten keine direkten Wahrscheinlichkeiten zu.
Stattdessen verwendet man eine Wahrscheinlichkeitsdichte (probability density function), die über Intervalle integriert wird.
\end{conceptbox}

Nach diesem Kapitel solltest du folgende Fragen beantworten können:

\begin{itemize}
    \item Was ist eine kontinuierliche Zufallsvariable (continuous random variable)?
    \item Warum ist \(p(X=x)\) bei kontinuierlichen Zufallsvariablen meistens \(0\)?
    \item Was ist eine Wahrscheinlichkeitsdichte (probability density function)?
    \item Wie berechnet man Wahrscheinlichkeiten über Intervalle?
    \item Welche Bedingungen muss eine Wahrscheinlichkeitsdichte erfüllen?
    \item Wie funktionieren Summenregel, Produktregel und Bayes-Regel für Dichten?
    \item Was ist ein Histogramm (histogram) und wie hängt es mit Dichteschätzung zusammen?
    \item Was sind Erwartungswert (expectation), Mittelwert (mean), Varianz (variance), Kovarianz (covariance) und Korrelation (correlation)?
    \item Was bedeutet Sampling-Approximation (sampling approximation)?
\end{itemize}

\section{Von diskreten zu kontinuierlichen Zufallsvariablen}

Eine diskrete Zufallsvariable \(X\) kann nur einzelne Werte annehmen.
Zum Beispiel beim Würfeln:

\[
    X \in \{1,2,3,4,5,6\}.
\]

Dann ist eine Wahrscheinlichkeit wie

\[
    p(X=3)
\]

sinnvoll und kann zum Beispiel den Wert

\[
    p(X=3) = \frac{1}{6}
\]

haben.

Bei einer kontinuierlichen Zufallsvariable ist die Situation anders.
Angenommen, \(X\) beschreibt eine Position auf einer Linie:

\[
    X \in \mathbb{R}.
\]

Dann gibt es unendlich viele mögliche Werte.
Die Wahrscheinlichkeit, genau einen einzelnen Wert zu treffen, ist in typischen kontinuierlichen Modellen gleich \(0\):

\[
    p(X=x) = 0.
\]

Das bedeutet aber nicht, dass nichts passieren kann.
Sinnvoll sind Wahrscheinlichkeiten für Intervalle:

\[
    p(X \in [a,b]).
\]

\begin{examplebox}
Wenn \(X\) die Körpergröße einer zufällig ausgewählten Person in Zentimetern beschreibt, ist die Wahrscheinlichkeit, exakt

\[
    X = 178{,}0000000000\ldots
\]

zu erhalten, praktisch \(0\).

Sinnvoll ist dagegen eine Frage wie:

\[
    p(175 \leq X \leq 180).
\]
\end{examplebox}

\begin{notebox}
Bei kontinuierlichen Zufallsvariablen sind einzelne Punkte normalerweise nicht wichtig.
Wichtig sind Bereiche, Intervalle oder allgemeiner Mengen.
\end{notebox}

\section{Wahrscheinlichkeitsdichten}

Für kontinuierliche Zufallsvariablen verwendet man Wahrscheinlichkeitsdichten (probability density functions).
Eine Wahrscheinlichkeitsdichte beschreibt nicht direkt die Wahrscheinlichkeit eines einzelnen Wertes, sondern wie stark Wahrscheinlichkeit in einem Bereich konzentriert ist.

\begin{definitionbox}
Eine Wahrscheinlichkeitsdichte (probability density function) einer kontinuierlichen Zufallsvariable \(X\) ist eine Funktion

\[
    p : \Omega_x \to \mathbb{R},
\]

wobei \(\Omega_x\) der Wertebereich von \(X\) ist.

Für eine eindimensionale kontinuierliche Zufallsvariable schreibt man meistens \(p(x)\).
\end{definitionbox}

Eine Dichte \(p(x)\) muss zwei wichtige Eigenschaften erfüllen:

\[
    p(x) \geq 0
    \quad
    \text{für alle}
    \quad
    x \in \Omega_x,
\]

und

\[
    \int_{x \in \Omega_x} p(x)\,dx = 1.
\]

Falls \(\Omega_x = \mathbb{R}\), schreibt man:

\[
    \int_{-\infty}^{\infty} p(x)\,dx = 1.
\]

\begin{conceptbox}
Eine Wahrscheinlichkeitsdichte muss nicht kleiner oder gleich \(1\) sein.

Wichtig ist nicht die Höhe der Dichte an einem Punkt, sondern die Fläche unter der Dichte.
Diese gesamte Fläche muss \(1\) ergeben.
\end{conceptbox}

\section{Wahrscheinlichkeiten als Integrale}

Bei kontinuierlichen Zufallsvariablen erhält man Wahrscheinlichkeiten durch Integration der Dichte.

\begin{definitionbox}
Für eine kontinuierliche Zufallsvariable \(X\) mit Dichte \(p(x)\) gilt:

\[
    p(X \in [a,b])
    =
    \int_a^b p(x)\,dx.
\]

Die Wahrscheinlichkeit, dass \(X\) im Intervall \([a,b]\) liegt, ist also die Fläche unter der Dichte zwischen \(a\) und \(b\).
\end{definitionbox}

\begin{examplebox}
Wenn \(p(x)\) die Dichte einer Messgröße ist, dann bedeutet

\[
    \int_2^5 p(x)\,dx
\]

die Wahrscheinlichkeit, dass der Messwert zwischen \(2\) und \(5\) liegt.
\end{examplebox}

\begin{notebox}
Für kontinuierliche Zufallsvariablen ist es meistens egal, ob man

\[
    [a,b],
    \quad
    (a,b),
    \quad
    [a,b),
    \quad
    (a,b]
\]

verwendet.
Die Randpunkte haben Wahrscheinlichkeit \(0\).
Daher gilt typischerweise:

\[
    p(X \in [a,b])
    =
    p(X \in (a,b)).
\]
\end{notebox}

\section{Die Gleichverteilung}

Ein einfaches Beispiel für eine kontinuierliche Verteilung ist die Gleichverteilung (uniform distribution).

\begin{definitionbox}
Eine kontinuierliche Zufallsvariable \(X\) ist gleichverteilt auf dem Intervall \([\alpha,\beta]\), wenn ihre Dichte gegeben ist durch

\[
    p(x \mid \alpha,\beta)
    =
    \begin{cases}
        \frac{1}{\beta-\alpha}, & x \in [\alpha,\beta], \\
        0, & \text{sonst}.
    \end{cases}
\]

Dabei gilt \(\alpha < \beta\).
\end{definitionbox}

Die Dichte ist auf dem Intervall konstant.
Außerhalb des Intervalls ist sie \(0\).

Die Normalisierung stimmt, denn:

\[
    \int_{-\infty}^{\infty} p(x \mid \alpha,\beta)\,dx
    =
    \int_{\alpha}^{\beta}
    \frac{1}{\beta-\alpha}\,dx.
\]

Da \(\frac{1}{\beta-\alpha}\) konstant ist, gilt:

\[
    \int_{\alpha}^{\beta}
    \frac{1}{\beta-\alpha}\,dx
    =
    \frac{1}{\beta-\alpha}
    \int_{\alpha}^{\beta} 1\,dx
    =
    \frac{1}{\beta-\alpha}
    (\beta-\alpha)
    =
    1.
\]

\begin{examplebox}
Für eine Gleichverteilung auf \([0,1]\) gilt:

\[
    p(x)
    =
    \begin{cases}
        1, & x \in [0,1], \\
        0, & \text{sonst}.
    \end{cases}
\]

Dann ist zum Beispiel:

\[
    p(X \in [0{,}2,0{,}7])
    =
    \int_{0{,}2}^{0{,}7} 1\,dx
    =
    0{,}5.
\]
\end{examplebox}

\section{Erwartungswert}

Der Erwartungswert (expectation) ist eine der wichtigsten Größen in der Wahrscheinlichkeitstheorie.
Er beschreibt den durchschnittlich erwarteten Wert einer Zufallsvariable oder einer Funktion einer Zufallsvariable.

\begin{definitionbox}
Sei \(X\) eine kontinuierliche Zufallsvariable mit Dichte \(p(x)\).
Für eine Funktion \(f(X)\) ist der Erwartungswert (expectation) definiert als

\[
    \langle f(X) \rangle_p
    =
    \int_{x \in \Omega_x}
    p(x) f(x)\,dx.
\]

Alternativ schreibt man oft:

\[
    \mathbb{E}_p[f(X)]
    =
    \int_{x \in \Omega_x}
    p(x) f(x)\,dx.
\]
\end{definitionbox}

Wenn klar ist, welche Verteilung gemeint ist, schreibt man auch einfach:

\[
    \mathbb{E}[f(X)].
\]

\begin{notebox}
Die Notationen

\[
    \langle f(X) \rangle_p
    \quad
    \text{und}
    \quad
    \mathbb{E}_p[f(X)]
\]

bedeuten dasselbe:
Erwartungswert von \(f(X)\) bezüglich der Verteilung \(p\).
\end{notebox}

\section{Mittelwert}

Der Mittelwert (mean) ist der Erwartungswert der Identitätsfunktion.
Das heißt, man setzt

\[
    f(X) = X.
\]

\begin{definitionbox}
Der Mittelwert (mean) einer kontinuierlichen Zufallsvariable \(X\) mit Dichte \(p(x)\) ist

\[
    \mu
    =
    \mathbb{E}[X]
    =
    \int_{x \in \Omega_x} p(x)x\,dx.
\]
\end{definitionbox}

Für \(\Omega_x = \mathbb{R}\) schreibt man:

\[
    \mu
    =
    \int_{-\infty}^{\infty} p(x)x\,dx.
\]

\begin{examplebox}
Für die Gleichverteilung auf \([\alpha,\beta]\) gilt:

\[
    p(x \mid \alpha,\beta)
    =
    \frac{1}{\beta-\alpha}
    \quad
    \text{für}
    \quad
    x \in [\alpha,\beta].
\]

Der Mittelwert ist:

\[
    \mathbb{E}[X]
    =
    \int_{\alpha}^{\beta}
    \frac{1}{\beta-\alpha}x\,dx.
\]

Also:

\[
    \mathbb{E}[X]
    =
    \frac{1}{\beta-\alpha}
    \int_{\alpha}^{\beta} x\,dx.
\]

Mit

\[
    \int_{\alpha}^{\beta} x\,dx
    =
    \left[
        \frac{x^2}{2}
    \right]_{\alpha}^{\beta}
    =
    \frac{\beta^2-\alpha^2}{2}
\]

folgt:

\[
    \mathbb{E}[X]
    =
    \frac{1}{\beta-\alpha}
    \frac{\beta^2-\alpha^2}{2}.
\]

Da

\[
    \beta^2-\alpha^2
    =
    (\beta-\alpha)(\beta+\alpha)
\]

gilt:

\[
    \mathbb{E}[X]
    =
    \frac{\alpha+\beta}{2}.
\]
\end{examplebox}

\begin{conceptbox}
Der Mittelwert der Gleichverteilung auf \([\alpha,\beta]\) liegt genau in der Mitte des Intervalls:

\[
    \mathbb{E}[X]
    =
    \frac{\alpha+\beta}{2}.
\]
\end{conceptbox}

\section{Varianz}

Der Mittelwert beschreibt die Lage einer Verteilung.
Er sagt aber nicht, wie stark die Werte um diesen Mittelwert streuen.
Dafür verwendet man die Varianz (variance).

\begin{definitionbox}
Die Varianz (variance) einer kontinuierlichen Zufallsvariable \(X\) mit Mittelwert \(\mu = \mathbb{E}[X]\) ist

\[
    \mathrm{Var}(X)
    =
    \mathbb{E}\left[(X-\mu)^2\right].
\]

Als Integral:

\[
    \mathrm{Var}(X)
    =
    \int_{x \in \Omega_x}
    p(x)(x-\mu)^2\,dx.
\]
\end{definitionbox}

Die Varianz ist immer nichtnegativ:

\[
    \mathrm{Var}(X) \geq 0.
\]

Das liegt daran, dass

\[
    (x-\mu)^2 \geq 0
\]

für alle \(x\) gilt.

\begin{conceptbox}
Die Varianz misst die mittlere quadratische Abweichung vom Mittelwert.

Große Varianz bedeutet:
Die Werte streuen stark.

Kleine Varianz bedeutet:
Die Werte liegen nah am Mittelwert.
\end{conceptbox}

\subsection{Alternative Formel für die Varianz}

Die Varianz kann auch anders geschrieben werden:

\[
    \mathrm{Var}(X)
    =
    \mathbb{E}[X^2]
    -
    \mathbb{E}[X]^2.
\]

Herleitung:

\[
    \mathrm{Var}(X)
    =
    \mathbb{E}\left[(X-\mu)^2\right].
\]

Ausmultiplizieren ergibt:

\[
    (X-\mu)^2
    =
    X^2 - 2\mu X + \mu^2.
\]

Also:

\[
    \mathrm{Var}(X)
    =
    \mathbb{E}[X^2 - 2\mu X + \mu^2].
\]

Wegen Linearität des Erwartungswerts gilt:

\[
    \mathrm{Var}(X)
    =
    \mathbb{E}[X^2]
    -
    2\mu \mathbb{E}[X]
    +
    \mathbb{E}[\mu^2].
\]

Da \(\mu\) konstant ist, gilt:

\[
    \mathbb{E}[\mu^2] = \mu^2.
\]

Außerdem ist

\[
    \mathbb{E}[X] = \mu.
\]

Damit:

\[
    \mathrm{Var}(X)
    =
    \mathbb{E}[X^2]
    -
    2\mu^2
    +
    \mu^2
    =
    \mathbb{E}[X^2]
    -
    \mu^2.
\]

Also:

\[
    \mathrm{Var}(X)
    =
    \mathbb{E}[X^2]
    -
    \mathbb{E}[X]^2.
\]

\begin{examtipbox}
Die Formel

\[
    \mathrm{Var}(X)
    =
    \mathbb{E}[X^2]
    -
    \mathbb{E}[X]^2
\]

ist oft praktisch für Rechnungen.
Sie ist aber nicht die Definition, sondern eine äquivalente Umformung.
\end{examtipbox}

\section{Standardabweichung}

Die Varianz hat quadrierte Einheiten.
Wenn \(X\) zum Beispiel in Metern gemessen wird, dann hat \(\mathrm{Var}(X)\) die Einheit Quadratmeter.
Deshalb verwendet man oft die Standardabweichung (standard deviation).

\begin{definitionbox}
Die Standardabweichung (standard deviation) ist die Wurzel der Varianz:

\[
    \sigma
    =
    \sqrt{\mathrm{Var}(X)}.
\]
\end{definitionbox}

Wenn man schreibt

\[
    \sigma^2 = \mathrm{Var}(X),
\]

dann ist \(\sigma\) die Standardabweichung und \(\sigma^2\) die Varianz.

\section{Kovarianz}

Bisher wurde die Streuung einer einzelnen Zufallsvariable betrachtet.
Oft interessiert man sich aber für den Zusammenhang zwischen zwei Zufallsvariablen.
Dafür verwendet man die Kovarianz (covariance).

\begin{definitionbox}
Seien \(X\) und \(Y\) zwei Zufallsvariablen mit Mittelwerten

\[
    \mu_X = \mathbb{E}[X],
    \qquad
    \mu_Y = \mathbb{E}[Y].
\]

Die Kovarianz (covariance) ist definiert als

\[
    \mathrm{Cov}(X,Y)
    =
    \mathbb{E}\left[(X-\mu_X)(Y-\mu_Y)\right].
\]
\end{definitionbox}

Für kontinuierliche Zufallsvariablen mit gemeinsamer Dichte \(p(x,y)\) gilt:

\[
    \mathrm{Cov}(X,Y)
    =
    \int
    \int
    p(x,y)
    (x-\mu_X)(y-\mu_Y)
    \,dx\,dy.
\]

\begin{conceptbox}
Die Kovarianz beschreibt, ob zwei Variablen gemeinsam steigen oder fallen.

\begin{itemize}
    \item \(\mathrm{Cov}(X,Y) > 0\): Große Werte von \(X\) treten eher mit großen Werten von \(Y\) auf.
    \item \(\mathrm{Cov}(X,Y) < 0\): Große Werte von \(X\) treten eher mit kleinen Werten von \(Y\) auf.
    \item \(\mathrm{Cov}(X,Y) \approx 0\): Kein klarer linearer Zusammenhang.
\end{itemize}
\end{conceptbox}

Eine alternative Formel ist:

\[
    \mathrm{Cov}(X,Y)
    =
    \mathbb{E}[XY]
    -
    \mathbb{E}[X]\mathbb{E}[Y].
\]

\section{Korrelation}

Die Kovarianz hängt von den Einheiten der Variablen ab.
Um einen einheitenlosen Zusammenhang zu erhalten, verwendet man die Korrelation (correlation).

\begin{definitionbox}
Die Korrelation (correlation) von \(X\) und \(Y\) ist

\[
    \rho_{X,Y}
    =
    \frac{
        \mathrm{Cov}(X,Y)
    }{
        \sigma_X \sigma_Y
    },
\]

wobei \(\sigma_X\) und \(\sigma_Y\) die Standardabweichungen von \(X\) und \(Y\) sind.
\end{definitionbox}

Es gilt:

\[
    -1 \leq \rho_{X,Y} \leq 1.
\]

\begin{itemize}
    \item \(\rho_{X,Y} = 1\): perfekter positiver linearer Zusammenhang.
    \item \(\rho_{X,Y} = -1\): perfekter negativer linearer Zusammenhang.
    \item \(\rho_{X,Y} = 0\): kein linearer Zusammenhang.
\end{itemize}

\begin{notebox}
Korrelation \(0\) bedeutet nicht automatisch Unabhängigkeit.
Sie bedeutet nur, dass kein linearer Zusammenhang gemessen wird.
Nichtlineare Zusammenhänge können trotzdem existieren.
\end{notebox}

\section{Gemeinsame Dichten}

Bei zwei kontinuierlichen Zufallsvariablen \(X\) und \(Y\) verwendet man eine gemeinsame Dichte (joint density)

\[
    p(x,y).
\]

\begin{definitionbox}
Eine gemeinsame Dichte (joint density) \(p(x,y)\) beschreibt die gemeinsame Verteilung zweier kontinuierlicher Zufallsvariablen \(X\) und \(Y\).

Sie muss erfüllen:

\[
    p(x,y) \geq 0
\]

und

\[
    \int_{-\infty}^{\infty}
    \int_{-\infty}^{\infty}
    p(x,y)\,dx\,dy
    =
    1.
\]
\end{definitionbox}

Die Wahrscheinlichkeit, dass \(X\) in \([a,b]\) und \(Y\) in \([c,d]\) liegt, ist:

\[
    p(X \in [a,b], Y \in [c,d])
    =
    \int_a^b
    \int_c^d
    p(x,y)\,dy\,dx.
\]

\begin{notebox}
Die Reihenfolge der Integrale ist hier nicht wesentlich, solange die Integrationsbereiche korrekt sind.
Je nach Konvention schreibt man auch

\[
    \int_c^d
    \int_a^b
    p(x,y)\,dx\,dy.
\]
\end{notebox}

\section{Marginalisierung bei Dichten}

Bei diskreten Zufallsvariablen erhält man Randwahrscheinlichkeiten durch Summieren.
Bei kontinuierlichen Zufallsvariablen ersetzt man Summen durch Integrale.

\begin{definitionbox}
Aus einer gemeinsamen Dichte \(p(x,y)\) erhält man die Randdichte (marginal density) von \(X\) durch Integration über \(y\):

\[
    p(x)
    =
    \int_{y \in \Omega_y}
    p(x,y)\,dy.
\]

Analog:

\[
    p(y)
    =
    \int_{x \in \Omega_x}
    p(x,y)\,dx.
\]
\end{definitionbox}

\begin{conceptbox}
Diskreter Fall:

\[
    p(x)
    =
    \sum_y p(x,y).
\]

Kontinuierlicher Fall:

\[
    p(x)
    =
    \int p(x,y)\,dy.
\]

Die Idee ist dieselbe:
Man entfernt eine nicht interessierende Variable durch Summieren oder Integrieren.
\end{conceptbox}

\section{Bedingte Dichten}

Auch für kontinuierliche Zufallsvariablen gibt es bedingte Dichten (conditional densities).

\begin{definitionbox}
Die bedingte Dichte (conditional density) von \(Y\) gegeben \(X=x\) ist

\[
    p(y \mid x)
    =
    \frac{p(x,y)}{p(x)}
\]

für \(p(x)>0\).
\end{definitionbox}

Daraus folgt die Produktregel für Dichten:

\[
    p(x,y)
    =
    p(y \mid x)p(x)
    =
    p(x \mid y)p(y).
\]

\begin{notebox}
Die Schreibweise sieht genauso aus wie im diskreten Fall.
Der Unterschied ist:
Bei kontinuierlichen Variablen sind \(p(x)\), \(p(y)\) und \(p(x,y)\) Dichten, keine Punktwahrscheinlichkeiten.
\end{notebox}

\section{Bayes-Regel für Dichten}

Die Bayes-Regel gilt auch für Wahrscheinlichkeitsdichten.
Für kontinuierliches \(x\) und diskretes \(y\) erhält man:

\[
    p(y \mid x)
    =
    \frac{
        p(x \mid y)p(y)
    }{
        \sum_{y'} p(x \mid y')p(y')
    }.
\]

Wenn auch \(y\) kontinuierlich ist, wird die Summe im Nenner durch ein Integral ersetzt:

\[
    p(y \mid x)
    =
    \frac{
        p(x \mid y)p(y)
    }{
        \int p(x \mid y')p(y')\,dy'
    }.
\]

\begin{definitionbox}
Bayes-Regel (Bayes' rule) für kontinuierliche Zufallsvariablen:

\[
    p(y \mid x)
    =
    \frac{
        p(x \mid y)p(y)
    }{
        p(x)
    },
\]

wobei

\[
    p(x)
    =
    \int p(x \mid y')p(y')\,dy'.
\]
\end{definitionbox}

\begin{examtipbox}
In kontinuierlichen Fällen gilt dieselbe Logik wie im diskreten Fall.

Der wichtigste Unterschied ist:

\[
    \sum
    \quad
    \text{wird ersetzt durch}
    \quad
    \int.
\]
\end{examtipbox}

\section{Histogramme und Dichteschätzung}

In der Praxis kennt man die echte Dichte \(p(x)\) oft nicht.
Man hat nur Datenpunkte

\[
    x_1,x_2,\dots,x_N.
\]

Eine einfache Methode zur Schätzung einer Dichte ist das Histogramm (histogram).

\begin{definitionbox}
Ein Histogramm (histogram) teilt den Wertebereich in Intervalle, sogenannte Bins (bins), und zählt, wie viele Datenpunkte in jedem Bin liegen.
\end{definitionbox}

Sei \(I_i\) ein Bin.
Die Breite des Bins sei

\[
    |I_i|.
\]

Die Anzahl der Datenpunkte im Bin \(I_i\) sei

\[
    n_i.
\]

Dann ist die relative Häufigkeit im Bin:

\[
    \frac{n_i}{N}.
\]

Diese relative Häufigkeit approximiert die Wahrscheinlichkeit, dass ein Datenpunkt in \(I_i\) liegt:

\[
    p(X \in I_i)
    \approx
    \frac{n_i}{N}.
\]

Andererseits gilt für die Dichte:

\[
    p(X \in I_i)
    =
    \int_{I_i} p(x)\,dx.
\]

Wenn das Intervall \(I_i\) klein ist und \(p(x)\) darin ungefähr konstant ist, dann gilt:

\[
    \int_{I_i} p(x)\,dx
    \approx
    p(\bar{x}_i)|I_i|,
\]

wobei \(\bar{x}_i\) der Mittelpunkt des Bins ist.
Damit:

\[
    \frac{n_i}{N}
    \approx
    p(\bar{x}_i)|I_i|.
\]

Also:

\[
    p(\bar{x}_i)
    \approx
    \frac{n_i}{N|I_i|}.
\]

\begin{conceptbox}
Die Histogramm-Schätzung einer Dichte lautet:

\[
    p(x)
    \approx
    \frac{n_i}{N|I_i|}
    \quad
    \text{für}
    \quad
    x \in I_i.
\]

Dabei ist \(n_i\) die Anzahl der Datenpunkte im Bin \(I_i\), \(N\) die Gesamtzahl der Datenpunkte und \(|I_i|\) die Binbreite.
\end{conceptbox}

\begin{notebox}
Wenn die Bins sehr breit sind, wird die Schätzung grob.
Wenn die Bins sehr schmal sind, wird die Schätzung stark verrauscht.
Die Wahl der Binbreite ist daher wichtig.
\end{notebox}

\section{Sampling-Approximation von Erwartungswerten}

Erwartungswerte sind Integrale.
Viele Integrale sind analytisch schwer oder gar nicht lösbar.
Wenn man aber Stichproben (samples)

\[
    x_1,x_2,\dots,x_N
\]

aus der Verteilung \(p(x)\) ziehen kann, kann man Erwartungswerte näherungsweise berechnen.

\begin{definitionbox}
Seien \(x_1,\dots,x_N\) Stichproben aus einer Verteilung \(p(x)\).
Dann kann man den Erwartungswert von \(f(X)\) approximieren durch

\[
    \mathbb{E}_p[f(X)]
    =
    \int p(x)f(x)\,dx
    \approx
    \frac{1}{N}
    \sum_{n=1}^{N} f(x_n).
\]

Das nennt man Sampling-Approximation (sampling approximation) oder Monte-Carlo-Approximation (Monte Carlo approximation).
\end{definitionbox}

\begin{conceptbox}
Die Sampling-Approximation ersetzt ein Integral durch einen Durchschnitt über Stichproben:

\[
    \int p(x)f(x)\,dx
    \approx
    \frac{1}{N}
    \sum_{n=1}^{N} f(x_n).
\]
\end{conceptbox}

Je größer \(N\) ist, desto besser wird die Approximation typischerweise.
Die Genauigkeit hängt aber auch von der Varianz der Funktion \(f(X)\) ab.

\section{Beispiel: Numerische Integration mit Sampling}

Wir betrachten das Integral

\[
    \int_{-\infty}^{\infty}
    e^{-x^2}x^2\,dx.
\]

Dieses Integral sieht nicht direkt wie ein Erwartungswert aus.
Wir können es aber so umformen, dass eine Normalverteilung erscheint.

Die Normalverteilung mit Mittelwert \(0\) und Varianz \(\frac{1}{2}\) hat die Dichte

\[
    \mathcal{N}\left(x \mid 0,\frac{1}{2}\right)
    =
    \frac{1}{\sqrt{\pi}}
    e^{-x^2}.
\]

Daraus folgt:

\[
    e^{-x^2}
    =
    \sqrt{\pi}
    \mathcal{N}\left(x \mid 0,\frac{1}{2}\right).
\]

Damit:

\[
    \int_{-\infty}^{\infty}
    e^{-x^2}x^2\,dx
    =
    \sqrt{\pi}
    \int_{-\infty}^{\infty}
    \mathcal{N}\left(x \mid 0,\frac{1}{2}\right)
    x^2\,dx.
\]

Das Integral auf der rechten Seite ist ein Erwartungswert:

\[
    \int_{-\infty}^{\infty}
    \mathcal{N}\left(x \mid 0,\frac{1}{2}\right)
    x^2\,dx
    =
    \mathbb{E}[X^2]
\]

für

\[
    X \sim \mathcal{N}\left(0,\frac{1}{2}\right).
\]

Mit Stichproben

\[
    x_n \sim \mathcal{N}\left(0,\frac{1}{2}\right)
\]

gilt daher:

\[
    \int_{-\infty}^{\infty}
    e^{-x^2}x^2\,dx
    \approx
    \frac{\sqrt{\pi}}{N}
    \sum_{n=1}^{N}
    x_n^2.
\]

\begin{examplebox}
Das Integral

\[
    \int_{-\infty}^{\infty}
    e^{-x^2}\cos(x)\,dx
\]

kann analog geschrieben werden als

\[
    \sqrt{\pi}
    \int_{-\infty}^{\infty}
    \mathcal{N}\left(x \mid 0,\frac{1}{2}\right)
    \cos(x)\,dx.
\]

Mit Stichproben \(x_n \sim \mathcal{N}(0,\frac{1}{2})\) erhält man:

\[
    \int_{-\infty}^{\infty}
    e^{-x^2}\cos(x)\,dx
    \approx
    \frac{\sqrt{\pi}}{N}
    \sum_{n=1}^{N}
    \cos(x_n).
\]
\end{examplebox}

\section{Analytische Lösung des ersten Integrals}

Das Integral

\[
    \int_{-\infty}^{\infty}
    e^{-x^2}x^2\,dx
\]

kann mit der Varianz der Normalverteilung interpretiert werden.

Für

\[
    X \sim \mathcal{N}\left(0,\frac{1}{2}\right)
\]

gilt:

\[
    \mathbb{E}[X] = 0
\]

und

\[
    \mathrm{Var}(X) = \frac{1}{2}.
\]

Da der Mittelwert \(0\) ist, gilt:

\[
    \mathrm{Var}(X)
    =
    \mathbb{E}[X^2].
\]

Also:

\[
    \mathbb{E}[X^2]
    =
    \frac{1}{2}.
\]

Daher:

\[
    \int_{-\infty}^{\infty}
    e^{-x^2}x^2\,dx
    =
    \sqrt{\pi}
    \mathbb{E}[X^2]
    =
    \sqrt{\pi}
    \frac{1}{2}.
\]

Somit:

\[
    \int_{-\infty}^{\infty}
    e^{-x^2}x^2\,dx
    =
    \frac{\sqrt{\pi}}{2}.
\]

\begin{examtipbox}
Die wichtige Idee ist nicht nur das Ergebnis, sondern die Umformung:

\[
    e^{-x^2}
    =
    \sqrt{\pi}
    \mathcal{N}\left(x \mid 0,\frac{1}{2}\right).
\]

Dadurch wird ein Integral zu einem Erwartungswert bezüglich einer bekannten Verteilung.
\end{examtipbox}

\section{Symmetrische Dichten}

Eine Dichte ist symmetrisch um \(\mu\), wenn links und rechts von \(\mu\) dieselbe Dichteform liegt.

\begin{definitionbox}
Eine Dichte \(p(x)\) ist symmetrisch um \(\mu\), wenn für alle \(x \in \mathbb{R}\) gilt:

\[
    p(\mu - x)
    =
    p(\mu + x).
\]
\end{definitionbox}

Für eine solche Dichte ist der Erwartungswert gleich dem Symmetriezentrum:

\[
    \mathbb{E}[X] = \mu.
\]

\subsection{Intuition}

Wenn die Dichte um \(\mu\) symmetrisch ist, dann gibt es zu jeder Abweichung nach rechts eine gleich wahrscheinliche Abweichung nach links.
Die Beiträge heben sich im Durchschnitt auf.
Deshalb liegt der Mittelwert genau bei \(\mu\).

\begin{examplebox}
Die Normalverteilung

\[
    \mathcal{N}(x \mid \mu,\sigma^2)
\]

ist symmetrisch um \(\mu\).
Daher gilt:

\[
    \mathbb{E}[X] = \mu.
\]
\end{examplebox}

\section{Wichtige Verteilungen}

In diesem Abschnitt sammeln wir einige wichtige kontinuierliche Verteilungen.
Sie werden später immer wieder verwendet.

\subsection{Gleichverteilung}

Die Gleichverteilung (uniform distribution) auf \([\alpha,\beta]\) ist:

\[
    p(x \mid \alpha,\beta)
    =
    \begin{cases}
        \frac{1}{\beta-\alpha}, & x \in [\alpha,\beta], \\
        0, & \text{sonst}.
    \end{cases}
\]

Der Mittelwert ist:

\[
    \mathbb{E}[X]
    =
    \frac{\alpha+\beta}{2}.
\]

Die Varianz ist:

\[
    \mathrm{Var}(X)
    =
    \frac{(\beta-\alpha)^2}{12}.
\]

\subsection{Normalverteilung}

Die Normalverteilung (Gaussian distribution) ist eine der wichtigsten Verteilungen im maschinellen Lernen.
In einer Dimension lautet ihre Dichte:

\[
    \mathcal{N}(x \mid \mu,\sigma^2)
    =
    \frac{1}{\sqrt{2\pi\sigma^2}}
    \exp\left(
        -\frac{(x-\mu)^2}{2\sigma^2}
    \right).
\]

Dabei ist \(\mu\) der Mittelwert und \(\sigma^2\) die Varianz.

\[
    \mathbb{E}[X] = \mu,
    \qquad
    \mathrm{Var}(X) = \sigma^2.
\]

\begin{conceptbox}
Die Normalverteilung ist wichtig, weil sie in vielen Modellen als Rauschmodell (noise model) verwendet wird.
Außerdem führt sie bei Maximum-Likelihood-Schätzung oft auf quadratische Fehlerfunktionen.
\end{conceptbox}

\section{Kontinuierliche Zufallsvektoren}

Bisher haben wir eindimensionale kontinuierliche Zufallsvariablen betrachtet.
In maschinellem Lernen sind Eingaben aber meistens Vektoren:

\[
    \vec{x}
    =
    \begin{pmatrix}
        x_1 \\
        \vdots \\
        x_D
    \end{pmatrix}
    \in \mathbb{R}^D.
\]

Dann spricht man von einem kontinuierlichen Zufallsvektor (continuous random vector).

\begin{definitionbox}
Ein kontinuierlicher Zufallsvektor \(\vec{X}\) mit Dichte \(p(\vec{x})\) erfüllt:

\[
    p(\vec{x}) \geq 0
    \quad
    \text{für alle}
    \quad
    \vec{x} \in \Omega_{\vec{x}},
\]

und

\[
    \int_{\vec{x} \in \Omega_{\vec{x}}}
    p(\vec{x})\,d\vec{x}
    =
    1.
\]
\end{definitionbox}

Für \(\vec{x} \in \mathbb{R}^D\) bedeutet \(d\vec{x}\):

\[
    d\vec{x}
    =
    dx_1\,dx_2\,\cdots\,dx_D.
\]

\begin{notebox}
In Machine Learning schreibt man häufig \(p(\vec{x})\) für die Dichte eines Datenvektors.
Das ist keine Wahrscheinlichkeit eines einzelnen Vektors, sondern eine Dichte.
Wahrscheinlichkeiten erhält man durch Integration über Bereiche im Datenraum.
\end{notebox}

\section{Zusammenhang zu maschinellem Lernen}

Kontinuierliche Zufallsvariablen treten im maschinellen Lernen an vielen Stellen auf:

\begin{itemize}
    \item Eingabevektoren \(\vec{x}\) werden oft als kontinuierlich modelliert.
    \item Regressionsziele \(t\) sind meistens kontinuierlich.
    \item Rauschen (noise) wird häufig mit Normalverteilungen modelliert.
    \item Parameter können in Bayesian Learning als Zufallsvariablen mit Dichten beschrieben werden.
    \item Gaussian Mixture Models verwenden Mischungen kontinuierlicher Dichten.
    \item PCA und probabilistische PCA arbeiten mit kontinuierlichen Vektoren.
\end{itemize}

\begin{conceptbox}
Viele ML-Modelle können als Aussagen über Dichten verstanden werden:

\[
    p(\vec{x}),
    \quad
    p(t \mid \vec{x}),
    \quad
    p(c \mid \vec{x}),
    \quad
    p(\vec{x} \mid c).
\]

Je nachdem, welche dieser Größen modelliert wird, entstehen unterschiedliche Arten von Lernverfahren.
\end{conceptbox}

\section{Häufige Fehler}

\begin{examtipbox}
Typische Fehler bei kontinuierlichen Zufallsvariablen:

\begin{enumerate}
    \item Eine Dichte \(p(x)\) mit einer Wahrscheinlichkeit \(p(X=x)\) verwechseln.
    \item Vergessen, dass Wahrscheinlichkeiten durch Integrale berechnet werden.
    \item Annehmen, dass eine Dichte nie größer als \(1\) sein darf.
    \item Die Normalisierungsbedingung \(\int p(x)\,dx=1\) vergessen.
    \item Bei Randdichten Summen statt Integrale verwenden.
    \item Bei Sampling-Approximation vergessen, dass die Stichproben aus der passenden Verteilung kommen müssen.
    \item Korrelation \(0\) mit Unabhängigkeit verwechseln.
\end{enumerate}
\end{examtipbox}

\section{Mini-Aufgaben}

\subsection{Aufgabe 1: Wahrscheinlichkeit bei Gleichverteilung}

Sei

\[
    X \sim \mathrm{Uniform}(0,4).
\]

Berechne

\[
    p(X \in [1,3]).
\]

\begin{examplebox}
Die Dichte ist

\[
    p(x)
    =
    \frac{1}{4}
    \quad
    \text{für}
    \quad
    x \in [0,4].
\]

Also:

\[
    p(X \in [1,3])
    =
    \int_1^3
    \frac{1}{4}\,dx
    =
    \frac{1}{4}
    (3-1)
    =
    \frac{1}{2}.
\]
\end{examplebox}

\subsection{Aufgabe 2: Mittelwert einer Gleichverteilung}

Sei

\[
    X \sim \mathrm{Uniform}(2,8).
\]

Berechne

\[
    \mathbb{E}[X].
\]

\begin{examplebox}
Für eine Gleichverteilung auf \([\alpha,\beta]\) gilt:

\[
    \mathbb{E}[X]
    =
    \frac{\alpha+\beta}{2}.
\]

Mit \(\alpha=2\) und \(\beta=8\):

\[
    \mathbb{E}[X]
    =
    \frac{2+8}{2}
    =
    5.
\]
\end{examplebox}

\subsection{Aufgabe 3: Varianz aus Momenten}

Angenommen,

\[
    \mathbb{E}[X] = 3
\]

und

\[
    \mathbb{E}[X^2] = 13.
\]

Berechne

\[
    \mathrm{Var}(X).
\]

\begin{examplebox}
Mit

\[
    \mathrm{Var}(X)
    =
    \mathbb{E}[X^2]
    -
    \mathbb{E}[X]^2
\]

folgt:

\[
    \mathrm{Var}(X)
    =
    13 - 3^2
    =
    13 - 9
    =
    4.
\]
\end{examplebox}

\section{Zusammenfassung}

\begin{itemize}
    \item Kontinuierliche Zufallsvariablen (continuous random variables) nehmen Werte in kontinuierlichen Räumen wie \(\mathbb{R}\) an.
    \item Einzelne Werte haben typischerweise Wahrscheinlichkeit \(0\).
    \item Wahrscheinlichkeiten über Intervalle berechnet man durch Integrale.
    \item Eine Wahrscheinlichkeitsdichte (probability density function) erfüllt:
    \[
        p(x) \geq 0
    \]
    und
    \[
        \int p(x)\,dx = 1.
    \]
    \item Für Intervalle gilt:
    \[
        p(X \in [a,b])
        =
        \int_a^b p(x)\,dx.
    \]
    \item Der Erwartungswert (expectation) ist:
    \[
        \mathbb{E}_p[f(X)]
        =
        \int p(x)f(x)\,dx.
    \]
    \item Der Mittelwert (mean) ist:
    \[
        \mathbb{E}[X]
        =
        \int p(x)x\,dx.
    \]
    \item Die Varianz (variance) ist:
    \[
        \mathrm{Var}(X)
        =
        \mathbb{E}\left[(X-\mu)^2\right].
    \]
    \item Die Kovarianz (covariance) misst den gemeinsamen linearen Zusammenhang zweier Variablen.
    \item Die Korrelation (correlation) ist die normierte Kovarianz.
    \item Summenregeln werden bei kontinuierlichen Variablen zu Integralen:
    \[
        p(x)
        =
        \int p(x,y)\,dy.
    \]
    \item Sampling-Approximation ersetzt Integrale durch Mittelwerte über Stichproben:
    \[
        \int p(x)f(x)\,dx
        \approx
        \frac{1}{N}
        \sum_{n=1}^{N}
        f(x_n).
    \]
\end{itemize}

\section{Typische Prüfungsfragen}

\begin{examtipbox}
Mögliche Prüfungsfragen zu diesem Kapitel:

\begin{enumerate}
    \item Warum ist \(p(X=x)=0\) für typische kontinuierliche Zufallsvariablen?
    \item Was ist eine Wahrscheinlichkeitsdichte (probability density function)?
    \item Welche zwei Bedingungen muss eine Dichte erfüllen?
    \item Wie berechnet man \(p(X \in [a,b])\)?
    \item Was ist der Unterschied zwischen Wahrscheinlichkeit und Dichte?
    \item Definiere Erwartungswert, Mittelwert und Varianz.
    \item Zeige, dass \(\mathrm{Var}(X)=\mathbb{E}[X^2]-\mathbb{E}[X]^2\).
    \item Was ist ein Histogramm und wie schätzt es eine Dichte?
    \item Erkläre die Sampling-Approximation eines Erwartungswerts.
    \item Wie werden Summenregel und Bayes-Regel im kontinuierlichen Fall geschrieben?
    \item Warum ist die Normalverteilung für ML wichtig?
    \item Was ist der Unterschied zwischen Kovarianz und Korrelation?
\end{enumerate}
\end{examtipbox}